\section{Complex series}
	\begin{itemize}
		\item \textit{Power series}: A series in powers of $(z-z_0)$ is known as a power series, i.e.
		\begin{align*}
			\sum_{n=0}^{\infty}{a_n(z-z_0)^n}&=a_0+a_1(z-z_0)+a_2(z-z_0)^2+\cdots+a_k(z-z_0)^k+\cdots
		\end{align*}
		where $z$ is a complex variable and $z_0$ is the centre of the series. The constants $a_0$, $a_1$, $a_2$, $\dots$ are known as the coefficients of the series.
		
		\item \textit{Region of convergence}: The region of convergence for a power series is a set of points over which the series is defined. There may be three distinct possibilities for the region of convergence:
		\begin{enumerate}[i.]
			\item The series converges only at the point $z=z_0$.
			\item The series converges for all points $z$ on the complex plane.
			\item The series converges everywhere inside a circular region given as $|z-z_0|<R$, where $z_0$ is the \textit{centre of convergence} and $R$ is its \textit{radius of convergence}.
		\end{enumerate}
		
		\item \textit{Radius of convergence}: Consider the power series which is centred at the origin, i.e. $\sum_{n=0}^{\infty}{a_nz^n}$. Then its radius of convergence $R$ is given by
		\begin{enumerate}[a)]
			\item $\dfrac{1}{R}=\lim_{n\to\infty}{\qty|a_n|^{1/n}}$
			
			\item $\dfrac{1}{R}=\lim_{n\to\infty}{\qty|\dfrac{a_{n+1}}{a_n}|}$
		\end{enumerate}
	\end{itemize}
	
	\subsection{Taylor's series expansion of a complex function}
	
	\subsection{Laurent's series expansion of a complex function}
	
	\subsubsection{Residue theorem}